\chapter*{Conclusion Générale}
\addcontentsline{toc}{chapter}{Conclusion Générale}
\markboth{Conclusion Générale}{Conclusion Générale}

\subsection*{Apport scientifique}
Ce travail contribue à l'état de l'art de l'IA appliquée à l'agriculture
par trois apports originaux. Premièrement, nous avons adapté et évalué
un pipeline RAG+LLM souverain pour la génération de texte en Darija
tunisienne — une langue dialectale non-standard très peu traitée dans
la littérature de NLP agricole — obtenant un Cohen's Kappa de 0.76 et
un BLEU-4 de 0.68 avec un modèle local de 7 milliards de paramètres.
Deuxièmement, nous avons proposé un schéma de validation statistique
rigoureux (validation croisée 5-fold avec TimeSeriesSplit, intervalles
de confiance à 95\%, analyse des menaces à la validité) pour les modèles
ML intégrés dans un ERP industriel, ce type de rigueur étant peu répandu
dans les publications sur les systèmes agricoles opérationnels. Troisièmement,
l'analyse comparative sur 3 runs indépendants YOLOv8 confirme la stabilité
du mAP@50 ($84.7\% \pm 0.45\%$) et constitue une donnée de référence pour
les travaux futurs sur la détection d'espèces animales méditerranéennes.

\subsection*{Apport technique}
Du point de vue de l'ingénierie logicielle, Smart Farm AI v3.0 Enterprise
démontre la faisabilité d'une plateforme agro-tech complète entièrement
construite sur des composants open-source (FastAPI, PostgreSQL, YOLOv8,
Ollama, ChromaDB, React). L'architecture multi-tenant avec fallback
3-niveaux (Ollama $\rightarrow$ Groq $\rightarrow$ statique), le déploiement Docker 5-services
avec HTTPS automatique Caddy et la PWA offline-first constituent un patron
d'architecture reproductible pour d'autres projets agri-tech en contexte
de connectivité limitée. Les 251 tests automatisés (98.8\% de taux de
réussite) et les benchmarks de performance (81 req/s, p50 = 94 ms) valident
la robustesse de la plateforme à l'échelle des fermes pilotes.

\subsection*{Apport sociétal}
L'impact sociétal de ce projet se manifeste à deux niveaux. Au niveau
individuel, les 7 fermes pilotes rapportent une réduction perçue de 23\%
des pertes animales et une satisfaction globale de 4.2/5, attestant d'une
adoption réelle et non seulement d'un intérêt de principe. Au niveau systémique,
Smart Farm AI représente un exemple de souveraineté numérique agricole :
en combinant coût nul (open-source), fonctionnement hors ligne et interface
en Darija tunisienne, il répond aux trois obstacles structurels à la
numérisation de l'agriculture familiale tunisienne identifiés dans l'analyse
de l'existant. Ce travail démontre que l'IA de pointe n'est pas une prérogative
des grandes exploitations industrielles ni des multinationales technologiques —
elle peut être déployée, évaluée rigoureusement et adoptée dans un contexte
rural méditerranéen avec des ressources matérielles et financières modestes.

% ============================================================
%  BIBLIOGRAPHIE
% ============================================================
% NOTE: \printbibliography and \appendix are in main.tex — do not repeat here.
